\scanpage{77}
\mspage{II. 25}
% archive sheet number: 76
% lines 1–4
\[(3.1.3)\]
\diagram{\begin{tikzcd}[ampersand replacement=\&,column sep=small]
1 \arrow[r,dash] \& \pi_1(\bar C,\xi_1) \arrow[r] \arrow[d,"\alpha"'] \& \pi_1(C,\xi_1) \arrow[r,dash] \arrow[d,"\beta"] \& \begin{gathered}\inserted{\operatorname{Gal}(\bar K/K)}\\\deleted{\unsure{\operatorname{Gal}(\bar K/K)}}\\\deleted{\unsure{\pi_1(\operatorname{Spec}K,\xi_1)}}\\\pi_1(\operatorname{Spec}K,\xi_1)\unsure{=}\end{gathered} \arrow[r,dash] \arrow[d,"\gamma"] \& 1\\
1 \arrow[r,dash] \& \pi_1(\bar X,\xi) \arrow[r,dash] \& \pi_1(X,\xi) \arrow[r,dash] \& \pi_1(\operatorname{Spec}k,\xi) \arrow[r,dash] \& 1
\end{tikzcd}}
,
% lines 5–6
où on a posé $\bar C=C\otimes_K\bar K$. \deleted{\unsure{Ceci prouve que}} Ceci précise le
% lines 7–9
\uline{Lemme} 3.2. Avec les \inserted{hyp. et} notations précédentes, dans le diagramme (3.1.3), les \inserted{trois} hom \deleted{$\gamma$ et} \inserted{verticaux} \deleted{\unsure{$\alpha$ sont des}} sont \uline{surjectifs}.
% lines 10–15
\deleted{\unsure{Dém.} Pour prouver le lemme}, Il suffit de prouver la surjectivité de $\gamma$ et $\alpha$. Celle de $\gamma$ provient du fait que $K=\unsure{k(t)}$ est une extension \deleted{transcendante} pure de $k$, donc \deleted{\unsure{le morphisme}} \deleted{\unsure{de}} $k$ est sép. fermé dans $K$.
% lines 16–23
Notons maintenant que la flèche $\bar C\longrightarrow\bar X$ donnant naissance à $\alpha$ est la même que celle qui \unsure{s'introduit|s'obtient} en \deleted{\unsure{remplaçant le cas}} faisant le changement de corps de base $k-\bar k$. Donc on peut supposer $k$ séparablement clos, et \deleted{\unsure{il est}} on \deleted{est ramené à \unsure{pr}} \inserted{doit examiner le morphisme sur les $\pi_1$ \uline{induit} par}
\[\unsure{\bar C|\underline C}-X.\]
% lines 24–26
Dire que le morphisme sur les $\pi_1$ est surjectif signifie que pour tt revêtement étale connexe $\bar X_1$ de $X$, son image inverse $\bar C_1$ sur $\bar C$ est
